\chapter{ニューラルネットワークの学習}

\section{はじめに}

【目的】

本章ではバックプロパゲーションを用いたニューラルネットワークの学習方法について理解します。

【目標】

\begin{itemize}
    \item バックプロパゲーションとは何か、について原理を説明できること。
    \item 最急降下法ならびに確率的勾配降下法とは何か、およびその原理について説明できること。
    \item 最急降下法を用いてニューラルネットワークのパラメータ更新の手順を自らの手で計算できること。
\end{itemize}

【章構成】

\begin{enumerate}
    \item ニューラルネットワークにおけるバックプロパゲーション
    \item 最急降下法を用いたニューラルネットワークのパラメータ最適化
    \item 例題1：最急降下法を用いた XOR 演算 を模した MLP の重みの更新
    \item 確率的勾配降下法とミニバッチ学習
\end{enumerate}

\section{ニューラルネットワークにおけるバックプロパゲーション}

\section{最急降下法を用いたパラメータ最適化}

\section{例題1：XOR演算MLPの重みの更新}

\section{確率的勾配降下法とミニバッチ学習}

本実験では扱わない。

\section{演習：MLPの学習機能}

前章で実装した MLP に学習機能を追加する。
前章で実装したのは MLP の出力の計算、つまり順伝播方向である。
本章では2つの機能を実装する。

\begin{enumerate}
    \item 逆伝播経路。
    \item 最急降下法を用いたパラメータの最適化。
\end{enumerate}

詳細は別紙に示す。
